\section{History records and complete symbolic domains}
\label{sec:bb-frontiers}

The balanced-stack theorem is complete.  The remaining realizations
use alternative recorders with tape data in two coordinates.
A motion of boxes is invertible.  To use it for an ordinary machine,
we must retain whatever information an instruction erases.  We use
finite history records in the sense of reversible computation, but
prove the local inverses below: reversibility of the initialized
execution would not by itself separate the full boxes.

A work instruction is $r=(q,a)\mapsto(p_r,b_r,d_r)$, with
$d_r\in\{-1,0,1\}$.  The work tape is indexed by $\Z$, starts with
head zero, and has only finitely many prescribed nonblank cells.
Halting states have no outgoing work instruction.  Missing instructions
can be replaced by write-preserving stationary steps to a fresh halt
state.  Unless a construction explicitly adds a dummy initial step,
an initially halting state is already a halt checkpoint. When a construction
uses a single halting state $h$, merge the original halting states and
redirect their incoming transitions; this preserves initial halting. In
the recorder tables $Q$ is the work-state set and $\Gamma$ the work
alphabet. Radix bases are chosen from the full cell alphabet, including
all history and mark tracks.

\subsection{A one-cell code and its exact inverse test}

A partial one-cell table sends a control and a scanned letter to a
new control, a written letter, and a displacement.  Its two incoming
conditions are: the new control determines the displacement, and the
new control together with the written letter determines the entire
old control and read letter.

\begin{lemma}[From one-cell tables to closed rectangles]
\label{lem:bb-onecell}\label{lem:bb-mid-prefix}
Give an alphabet of size $m$ digits separated by at least two, between
$0$ and $B-2$.  Put
\[
E(A)=\sum_{j\ge1}d(A_j)B^{-j},\qquad
I(v)=\sum_{j=1}^{|v|}d(v_j)B^{-j}+[0,B^{-|v|}].
\]
For a head-relative tape $(a_j)_{j\in\Z}$ use the coordinates
\[
 (x,y)=\big(E(a_{-1}a_{-2}\cdots),E(a_0a_1\cdots)\big).
\]
After writing at cell zero, a displacement $e$ makes the old cell $e$
the new cell zero.
A deterministic partial one-cell table satisfying the incoming
conditions acts, after separate state translations, by a finite list
of reciprocal affine maps on full closed rectangles.  Each source
family and each target family is separately disjoint.  For a rule
reading $a$, writing $b$, and moving by $e$, the branches are
\begin{equation}\label{eq:bb-one-cell-branches}
\begin{array}{c|c|c|c}
e&\text{source}&\text{target}&(x',y')\\\hline
1 &[0,1]\times I(a)&I(b)\times[0,1]
 &((d(b)+x)/B,By-d(a))\\
0 &[0,1]\times I(a)&[0,1]\times I(b)
 &(x,y+(d(b)-d(a))/B)\\
-1&I(c)\times I(a)&[0,1]\times(d(c)+I(b))/B
 &(Bx-d(c),[d(c)+(d(b)+By-d(a))/B]/B).
\end{array}
\end{equation}
The last row is split over every letter $c$.  Each linear part is
$\operatorname{diag}(B^{-e},B^e)$.  With odd digits
$1,3,\ldots,2m-1$ and $B\in\{2m+1,2m+2\}$, every code lies in
the interior of its applicable rectangle, and unscaled within-family
gaps are at least $B^{-2}$.
\end{lemma}
\begin{proof}
The identity $E(aA)=(d(a)+E(A))/B$ gives all three formulas and their
action on the independent tails.  At the first differing digit,
incompatible prefix intervals have a positive gap: the digit difference
is at least two while the child interval has length one in that scale.
Thus the source state, read letter, and any left split separate sources.
At one target state the displacement is fixed.  For a right or stationary
move the written letter separates the relevant target coordinate;
for a left move the pair $(c,b)$ separates the two-letter prefix.
The incoming inverse condition rules out duplicates.  These arguments
concern the complete intervals, including their endpoints.
For odd digits the infinite tail lies between $1/(B-1)$ and
$(2m-1)/(B-1)<1$, proving interiority.  The longest target prefix has
length two, giving the asserted gap.  Uniformly scaling state squares
by $s$ scales gaps by $s$ and leaves the linear parts unchanged.
\end{proof}

\subsection{Moving the work head before recording}

We first place a frontier to the right of every possible work head.
A temporary mark lets the history excursion return to the correct
cell even though its length grows with the computation.

\begin{lemma}[Move-first recorder]
\label{lem:bb-movefirst}\label{lem:bb-mid-right-compiler}
There is an effective one-cell table satisfying the incoming conditions
whose checkpoint configurations simulate the work machine.  After
$n$ work steps its frontier is at $n+2$, its records occupy
$2,\ldots,n+1$, and its temporary marks are all zero.  Each nonhalting
checkpoint has a finite positive continuation to the next one, with
no intervening checkpoint.
\end{lemma}
\begin{proof}
Use letters $(c,\ell,j)$ with work letter $c$, history
$\ell\in\{e,F\}\sqcup\{r\}$ and mark $j\in\{0,1\}$.
The controls are $C_q,J_r,S_r,P_p,T_p$.  The complete partial table is
\[
\begin{array}{lll}
 C_q:(a,\ell,0)&\mapsto J_r:(b_r,\ell,0),d_r&\ell\ne F,\\
 J_r:(c,\ell,0)&\mapsto S_r:(c,\ell,1),1&\ell\ne F,\\
 S_r:(c,\ell,0)&\mapsto S_r:(c,\ell,0),1&\ell\ne F,\\
 S_r:(c,F,0)&\mapsto P_{p_r}:(c,r,0),1,&\\
 P_p:(c,e,0)&\mapsto T_p:(c,F,0),-1,&\\
 T_p:(c,\ell,0)&\mapsto T_p:(c,\ell,0),-1&\ell\ne F,\\
 T_p:(c,\ell,1)&\mapsto C_p:(c,\ell,0),0&\ell\ne F.
\end{array}
\]
Initialize the frontier at two, all other history empty, and all marks
zero.  At checkpoint $n$ the head satisfies $|h_n|\le n$.
The first row updates the work tape and moves to $h'\le n+1$, strictly
left of the frontier $f=n+2$.  The second row marks $h'$.  The forward
scan writes $r$ at $f$, puts a new frontier at $f+1$, and the backward
scan stops exactly at the mark, which it clears.  Its length is
$2(f-h')+4$, so it is finite and positive.  This proves the invariant
inductively, including the exact work update and halt equivalence.

For the full-domain inverse, $r$ recovers the old work pair into $J_r$;
mark one distinguishes entry into $S_r$ from its mark-zero loop;
a written record determines entry into $P_p$; a written $F$
distinguishes entry into $T_p$ from its loop; and entry into $C_p$
uniquely restores mark one.  Every destination has the single incoming
displacement displayed in the table.  No invariant on remote cells
was used.
\end{proof}

\subsection{Recording before the work-head displacement}

The next table leaves its mark at the old work head.  Its final step
performs the work displacement; storing that displacement in the
checkpoint control keeps incoming moves unambiguous.

\begin{lemma}[Direction-recording scan]
\label{lem:bb-direction}\label{lem:bb-mid-mark-compiler}
Let $D=\{-1,0,1\}$ and let $O$ be either a singleton or $\{0,1\}$.
An effective one-cell table on
$\Gamma\times(\{\bot,F\}\sqcup\mathcal R)\times\{0,1\}\times O$,
where $\mathcal R=\{(q,e,a):q\text{ nonhalting},e\in D,a\in\Gamma\}$,
satisfies the incoming conditions.  It preserves the $O$ track in
its physical cells.  Starting with frontier at one, its checkpoint
at work step $n$ has frontier $n+1$, records at $1,\ldots,n$, zero
marks, and the correct work configuration.  A step from head $i$
takes $5+2(n-i)$ local rules.
\end{lemma}
\begin{proof}\label{proof:bb-mid-mark-compiler}
Use main controls $M(q,e)$, forward controls $G(r)$, and $A(p,d),B(p,d)$
for every work state $p$ and $d\in D$, including pairs not produced by
a work instruction. Copy the $O$ component in each row. In the first
three rows use every $r=(q,e,a)\in\mathcal R$ with
$\delta(q,a)=(p,c,d)$. The final three rows apply independently to
every $(p,d)\in Q\times D$, including pairs that are not outputs of
any work instruction. The variables $v\in\Gamma$ and
$h\in\{\bot,F\}\sqcup\mathcal R$ range over all letters subject to
the displayed restrictions. The whole table is
\[
\begin{array}{lll}
 M(q,e):(a,h,0)&\mapsto G(r):(c,h,1),1&h\ne F,\\
 G(r):(v,h,0)&\mapsto G(r):(v,h,0),1&h\ne F,\\
 G(r):(v,F,0)&\mapsto A(p,d):(v,r,0),1,&\\
 A(p,d):(v,\bot,0)&\mapsto B(p,d):(v,F,0),-1,&\\
 B(p,d):(v,h,0)&\mapsto B(p,d):(v,h,0),-1&h\ne F,\\
 B(p,d):(v,h,1)&\mapsto M(p,d):(v,h,0),d&h\ne F.
\end{array}
\]
Starting at $M(q_0,0)$, the first row writes and marks the old head.
The scan reaches frontier $j=n+1$, stores its record, creates the
next frontier, returns to the unique mark and makes the work move.
The count is $2(j-i)+3=5+2(n-i)$.  Since $i\le n<j$, every test is
satisfied and both scans are finite.

Into $G(r)$ the written mark distinguishes entry and loop; $r$
recovers the old state and work symbol.  Into $A(p,d)$ the record
recovers the old $G(r)$.  Into $B(p,d)$ a written frontier distinguishes
arrival and loop.  Into $M(p,d)$ only clearing a mark is possible.
All remaining components were copied.  This proves the inverse on
arbitrary tapes and the fixed incoming displacement.  Deleting the
passive $O$ coordinate is an exact specialization of every rule,
not merely an equivalence of initialized runs.  When $O=\{0,1\}$,
a unique initial one at physical cell zero records the absolute origin.
\end{proof}

\subsection{A frontier on the left and post-halt continuation}

\begin{lemma}[Two left-frontier tables]\label{lem:bb-mid-left-compiler}
There are two effective incoming-separated tables: a stopping table,
and a table which continues with idle work steps after halting.
At checkpoint $n$ their frontier is $f_n=-2-n$ and their marks vanish.
If the next work move ends at $i'$, the next checkpoint is reached
after $4+2(i'-f_n)$ local rules.
\end{lemma}
\begin{proof}
For the stopping version let $J=(Q\setminus\{h\})\times\Gamma$.
For the continuing version take $J=Q\times\Gamma$ and add
$\delta(h,a)=(h,a,0)$.  Use letters
$(c,\eta,j)\in\Gamma\times(\{\bot,F\}\sqcup J)\times\{0,1\}$,
controls $C_q,D_q,T_q,P_r,S_r$, and precisely
\begin{equation}\label{eq:bb-mid-left-table}
\begin{aligned}
 C_q:(a,\eta,0)&\mapsto P_r:(b_r,\eta,0),d_r,\\
 P_r:(c,\eta,0)&\mapsto S_r:(c,\eta,1),-1&&\eta\ne F,\\
 S_r:(c,\eta,0)&\mapsto S_r:(c,\eta,0),-1&&\eta\ne F,\\
 S_r:(c,F,0)&\mapsto D_{p_r}:(c,r,0),-1,\\
 D_q:(c,\bot,0)&\mapsto T_q:(c,F,0),1,\\
 T_q:(c,\eta,0)&\mapsto T_q:(c,\eta,0),1&&\eta\ne F,\\
 T_q:(c,\eta,1)&\mapsto C_q:(c,\eta,0),0&&\eta\ne F.
\end{aligned}
\end{equation}
In the first row only, the stopping version requires $\eta\ne F$;
the continuing version allows every $\eta$.  These domains remain
distinct even though initialized runs never use that difference.

Initialize $F$ at $-2$, empty history elsewhere and zero marks.
At checkpoint $n$, $i\ge-n=f_n+2$.  After the work update,
$i'\ge f_n+1$.  Marking $i'$ and scanning left therefore reaches
$f_n$, records $r$, puts the next frontier at $f_n-1$, and returns
right to the mark without crossing the new frontier.  Clearing it
completes the next work checkpoint.  The two finite traversals and
four endpoint actions give the stated count.  The continuing version
has a next rule even after every halt checkpoint.

Into $P_r$ the record index recovers the work pair.  Into $S_r$
mark one distinguishes entry from loop; into $D_q$ the written record
recovers the scan; into $T_q$ a written $F$ distinguishes entry from
loop; into $C_q$ there is only stationary unmarking.  This proves the
incoming conditions for both full partial domains.  It also gives
the configuration inverse: reverse the displacement, inspect the
written cell, and invert its uniquely determined rule.
\end{proof}

\subsection{A guarded update on two neighboring history cells}

One may advance the frontier and write its record in a single local
instruction.  The resulting return loop needs a neighboring-cell guard.
Without that guard, its image would overlap the logging image.

\begin{lemma}[Three-cell guarded window]\label{lem:bb-window-compiler}
There is a finite partial injective map on relative tapes with controls
$\mathrm{Ready}(q,e)$, $\mathrm{Go}(r)$ and $\mathrm{Back}(q,d)$.
Here $r=(q,e,a)$ records an instruction, with
$\delta(q,a)=(p(r),b(r),d(r))$.
Use work, history and mark tracks $A,H,J$.  At cursor position $i$,
the Ready row has $a=A_i$ and $r=(q,e,a)$.  The entire table is
\[
\begin{array}{c|l|l}
\text{control}&\text{test}&\text{action}\\\hline
\mathrm{Ready}(q,e)&q\text{ nonhalting},J_i=0
 &A_i\gets b(r),\ J_i\gets1,\ \mathrm{Go}(r),+1\\
\mathrm{Go}(r)&J_i=0,H_i\ne P&\mathrm{Go}(r),+1\\
\mathrm{Go}(r)&J_i=0,H_i=P,H_{i+1}=E
 &H_i\gets r,H_{i+1}\gets P,\ \mathrm{Back}(p(r),d(r)),-1\\
\mathrm{Back}(q,d)&J_i=0,H_{i+1}\ne P&\mathrm{Back}(q,d),-1\\
\mathrm{Back}(q,d)&J_i=1&J_i\gets0,\ \mathrm{Ready}(q,d),d.
\end{array}
\]
Unmentioned cells and tracks are copied.  Initialized with frontier
at $r_0\in\{1,2\}$, other history $E$, zero marks and
$\mathrm{Ready}(q_0,0)$, the checkpoint at step $n$ has frontier
$p_n=n+r_0$ and correct head
$h_n$.  Its next work step uses exactly
$2(p_n-h_n)+1\le4n+2r_0+1$ local rules.
\end{lemma}
\begin{proof}
The first rule writes and marks $h_n$.  Since $p_n>h_n$,
the forward scan reaches the frontier, moves it one cell right, and
the return scan reaches the mark.  Its neighboring-cell guard holds
because the new frontier is $p_n+1$.  The last rule clears the mark
and performs the work displacement.  Counting the two traversals and
endpoint actions gives $2(p_n-h_n)+1$; induction gives $|h_n|\le n$.

For injectivity on all relative tapes, into Go the old cursor is
output offset $-1$: a mark there distinguishes entry from the loop,
and $r$ restores the overwritten work letter and old Ready control.
Into Back the old cursor is offset $+1$.  An output pointer at $+2$
identifies a log insertion, with record at $+1$ recovering $r$ and
old pair $(P,E)$.  A nonpointer at $+2$ identifies the guarded loop.
Into Ready the corresponding Back control uniquely restores the
cleared mark.  These are all possibilities, proving the full inverse.
\end{proof}

\begin{lemma}[Closed rectangles for a finite guarded window]
\label{lem:bb-window-rectangles}
Suppose a finite partial injective local tape map reads and writes only
indices $[-a,b-1]$, then shifts the cursor by $e$, where
$-a\le e\le b$.  Split its domain by every complete word in that
window.  With gapped radix coding the input prefix lengths are
$(a,b)$ and output lengths $(a+e,b-e)$; the affine branch has linear
part $\operatorname{diag}(B^{-e},B^e)$.  Both full rectangle families
are separately disjoint.
\end{lemma}
\begin{proof}
After the writes the new left prefix consists of old indices
$e-1,e-2,\ldots,-a$, and the right prefix of $e,e+1,\ldots,b-1$.
The two remaining tails are free and unchanged, so the image is the
whole output cylinder.  Distinct source cylinders have disjoint
images by injectivity.  If both pairs of output prefixes were
compatible, a common extension would belong to both images; hence
at least one coordinate has incompatible prefixes.  The gapped
interval argument then separates the complete closed rectangles.
The prefix lengths give the two reciprocal scales directly.
\end{proof}
